\section{Experiments}
\label{sec:experiments}

\label{sec:results}

\subsection{Setup}

\label{sec:setup}

\textbf{Models and agent.} We evaluate DeepSeek V3.1, Kimi K2.5, Claude Sonnet 4,
Claude Sonnet 4.5, and GLM-5 through Amazon Bedrock at temperature 0. Each
model runs as a ReAct-style agent against the environment's MCP tool servers.

\textbf{Runs.} Each model--defense configuration runs all 90 retained cases without an injection, with the compound attack, and with every retained single-primitive attack: 90 benign and 1,433 attacked runs (90 compound, 1,343 single-primitive).

\textbf{Defenses.} We compare seven defenses spanning five mechanism families.
Spotlighting delimits tool output \citep{hines2024spotlighting}, while Sandwich
repeats the user request after each result \citep{debenedetti2024agentdojo}.
DRIFT-iso applies DRIFT's
injection-removal isolator; DRIFT-full additionally checks calls against a
planned tool sequence and parameter checklist \citep{li2025drift}. Progent enforces
LLM-generated JSON Schema policies and permits runtime updates from tool
results \citep{shi2025progent}. CaMeL tracks value provenance and requires control
arguments, such as a recipient or amount, to come from trusted fields; we write
its policies for our actions and define trust labels at the field level of each
tool result \citep{debenedetti2025defeating}. ARGUS releases a state-changing
call only if the request or content it does not detect as anomalous supports the
call, and the call violates no request-derived task invariant
\citep{weng2026argus}.

Defense-internal LLM calls use the agent
model.

\textbf{Controlled variation.} We compare each original single-primitive attack
with a version rewritten for up to three rounds against a carrier-only probe,
and with two expanded versions that keep the original attack verbatim and add
either plain or detection-optimized support. The probe uses DeepSeek V3.2 at
temperature 0 and receives only the \textbf{carrier}, the item that contains the attack
(such as an email or review), reconstructed from the tool output, but not the
user request. It flags instructions or other manipulative insertions that do
not fit the carrier, including attempts to influence actions or decisions.
Detection-optimized rewrites must remain valid, target only their primitive,
and alone justify the harmful outcome. Added support may span primitives, but
an ablation check requires that it cannot justify the harmful outcome on its
own. Rewriting optimizes evasion of the probe, not attack success.

We evaluate all four conditions without a defense and under
DRIFT-iso for
Claude Sonnet 4.5, DeepSeek V3.1, Kimi K2.5, and GLM-5.

\textbf{Trajectory annotation.} As a secondary behavioral analysis, Claude Opus
4.6 labels trajectories of the first single-primitive attack per case and
primitive for GLM-5 and Claude Sonnet 4.5 (up to
442 per model--configuration).
The judge receives the user request, planted content, expected sink, full
execution trace, and defense events. It records whether the content was
\emph{followed}, explicitly \emph{resisted}, present with \emph{no observable engagement},
\emph{absent} from the trace, or whether the trajectory is \emph{incomplete}, and the grounds for any resistance. Runtime interventions are measured
separately: content removal from the result delivered to the agent, and action
blocking from recorded policy decisions on emitted tool calls. These annotations
explain how outcomes arise; they do not determine ASR or TSR.

\textbf{Metrics.} ASR is the fraction of attacked runs in which the predefined
harmful deviation occurs; benign and attacked TSR measure task completion
without and with planted content (\S\ref{sec:evaluation}). Model-level ASR
pools single-primitive and compound runs, while primitive-level ASR uses
only single-primitive runs. Means are unweighted across models. Outcome tables
use all retained attacks; trajectory analyses use the first attack per case
and primitive.
